\section{Prüfer 双射与有标号树计数}\label{knowledge-prufer}\index{Prüfer sequence}\index{Cayley formula}
本节的树是顶点标号固定为 $1,\ldots,n$ 的简单无向树，不把同构的不同标号树合并。
先设 $n\ge2$。反复删除当前\textbf{最小标号叶子}，记录它唯一的邻点，直到剩两个点；
得到长度 $n-2$、每项属于 $[n]$ 的 Prüfer 序列。
$n=2$ 对应唯一空序列；$n=1$ 另有唯一空边树，度数为0，不代入含 $n-2$ 的公式。

\subsection{双射为何成立}
解码时保留所有尚未删去的标号。读下一项 $u$，
取当前未删除且没有出现在\textbf{整个剩余序列}中的最小标号 $v$，连边 $uv$，
删除 $v$ 并消耗这一项 $u$；最后连接剩下的两个标号。
当前有 $t$ 个点、$t-2$ 个剩余序列项，所以至少两个标号没有出现，可选点总存在。
$u$ 自己仍在序列中，故不会选出自环。每个被删点只连接到更晚删去的点，
最终得到一棵连通的 $n-1$ 边树。

关键不变式是：在当前剩余树中，标号 $i$ 在剩余序列里的出现次数等于 $d_i-1$。
每个点最后一条关联边在它作为叶子被删除时移除，或在最后两点之间保留到结束；这一条不记录该点。
此前每删除一个邻点，就记录一次它。
每一步删除一个叶子，会同时使邻点的度数和出现次数减少1，因此不变式继续成立。
反过来看解码得到的剩余树，每个活跃点有一条最终连接边，再加上剩余序列每次出现它时接来的边；
其度数同样是剩余出现次数加1。
所以未出现的活跃标号恰好是叶子；两边都取最小者，编码和解码逐步互为逆操作。
特别地，原树满足
\[
 d_i=1+\#\{\text{序列中等于 }i\text{ 的项}\}.
\]
本节给出计数依据，不新增可调用的编码/解码 API。

\subsection{Cayley、指定度数与叶子数}
任意长度 $n-2$ 的序列都对应唯一树，故树总数为 $n^{n-2}$。
给定整数度数 $d_1,\ldots,d_n$，若某个 $d_i<1$ 或 $\sum_i d_i\ne2n-2$，答案为0；
否则各标号出现次数已确定，答案是多项式系数
\[
 \frac{(n-2)!}{\prod_{i=1}^n(d_i-1)!}.
\]
例如 $n=4$，度数 $(3,1,1,1)$ 给出唯一指定中心的星形树，
度数 $(2,2,1,1)$ 则有2棵，而不是把两名度为2的顶点视作不可区分。

叶子恰是没有出现的标号。恰有 $\ell$ 个叶子的树数为
\[
 \binom n\ell\,(n-\ell)!\,S(n-2,n-\ell),
\]
其中 $S$ 为第二类 Stirling 数：先选叶子，再把序列位置满射到其余标号。
$n\ge3$ 时合法范围为 $2\le\ell\le n-1$；其他叶子数为0。
$n=2,\ell=2$ 时用 $S(0,0)=1$ 得1。若把叶子定义为度数1，单点树没有叶子，仍单独处理。
这里是整数计数恒等式；取模时不能无条件把阶乘分母换成逆元。
素数模的预处理与下标限制见第~\ref{compact-Binomial}~节（第~\pageref{compact-Binomial}~页），
一般模数条件见数学分册“模运算的适用条件”。
有根标号树及其生成函数见数学分册“隐式生成函数与拉格朗日反演”，不重复另造树计数模板。

\clearpage
\subsection{加权 Cayley 与连接已有连通块}
若完整图的边权有特殊乘积形式 $w_{ij}=x_i x_j$，则
\[
 \sum_T\prod_{ij\in T}w_{ij}
 =\sum_T\prod_i x_i^{d(i)}
 =\left(\prod_i x_i\right)\left(\sum_i x_i\right)^{n-2}.
\]
这里 $d(i)$ 表示点 $i$ 的度数。第二步利用出现次数恒等式，把每棵树变成公共因子 $\prod_i x_i$ 乘以其序列各项的权；
对所有序列求和就是展开右侧的幂。
证明没有除以任何 $x_i$，所以零权、负权以及任意正模数都可按多项式恒等式使用。
$n=2$ 的零次幂是常数1，即使 $\sum_i x_i=0$ 也如此。
这只适用于上述乘积边权；缺边图或一般边权应使用矩阵树定理，
见第~\ref{knowledge-matrix-tree-weighted}~节（第~\pageref{knowledge-matrix-tree-weighted}~页）及第~\ref{compact-MatrixTreeMod}~节（第~\pageref{compact-MatrixTreeMod}~页）。
代数权值和为零也可能来自抵消，不能据此断言没有树。

设一个固定图有 $k\ge2$ 个连通块，大小为正整数 $s_1,\ldots,s_k$，总顶点数 $N=\sum_i s_i$。
允许任选不同连通块的顶点对加边，使全图连通，且恰加 $k-1$ 条边。
缩点后新增边必须组成树；一条块间边 $ij$ 有 $s_i s_j$ 种端点选择。
因此不同\textbf{无序新增边集}的数量是
\[
 N^{k-2}\prod_{i=1}^k s_i.
\]
$k=1$ 时无需加边，答案为1，不计算负指数。原图的内部边固定，不计内部生成树选择。
若禁止部分端点对、给每条新增边限制，或数加边的执行顺序，就不能直接使用此式。
有标号树也不能简单除以 $n!$ 得到无标号树数，因为不同树的自同构数量不同。

双射背景见
\href{https://math.mit.edu/~rstan/algcomb/algcomb.pdf}{Stanley, Topics in Algebraic Combinatorics}，
导航与基础公式见
\href{https://github.com/OI-wiki/OI-wiki/blob/bc070e827180fbd75c2e27a16c1212f1d671d949/docs/graph/prufer.md}{固定版本 OI Wiki Prüfer 序列}。
以上用最小叶子的一致约定推导；不将本节计数说明当作整个上游页面的实现覆盖。
